% ============================================================
% abstract.tex
% ============================================================

\begin{abstract}
\noindent
Understanding how behavioral traits relate to cortical morphometry, and how those relationships shift with age, is a central question in human neuroscience.
We present a dual-cohort analysis that examines sleep quality, fluid intelligence, and neuroticism as predictors of cortical thickness and surface area, alongside a characterization of age-related cortical decline.
Using the Human Connectome Project Young Adult dataset (\textit{N}\,=\,1{,}104; ages 22--37; Desikan-Killiany atlas, 68 regions of interest) and the Aging Brain and Behaviour Cohort (\textit{N}\,=\,1{,}150; ages 36--89; Glasser HCP-MMP1 atlas, 360 regions of interest), we conducted region-wise Welch \textit{t}-tests, Pearson correlations, and one-way ANOVAs, applying Benjamini-Hochberg false discovery rate correction throughout.
Three principal findings emerged: (1)~fluid intelligence correlated positively with surface area across all 68 regions in young adults; (2)~cortical thickness and volume declined near-universally with age, with the strongest effects in primary motor cortex (\textit{r}\,=\,$-$0.69); and (3)~neuroticism was associated with widespread cortical thinning in the aging cohort (327 of 360 regions) but not in young adults, suggesting that cumulative stress exposure may be required for structural differences to manifest.
These results underscore the importance of cohort-specific designs when investigating brain-behavior associations and provide converging evidence that the aging cortex is differentially vulnerable to personality-linked physiological processes.
\end{abstract}

\vspace{0.5cm}
\noindent\textbf{Keywords:} cortical thickness, surface area, voxel-based morphometry, aging, sleep quality, fluid intelligence, neuroticism, Human Connectome Project
\newpage
